% 7 Threats to validity and open problems -- 0.6 page
\section{Threats to Validity and Open Problems}
\label{sec:threats}

\textbf{Internal validity.} We developed \sysname against the bench, so it may fit the bench better than unseen problems. Three things limit this: the answer keys come from the programs' fault switches, not from \sysname; the question thresholds were set by measuring every trace before \sysname's final runs; and the configuration used here was frozen and fingerprinted before them. Still, the fixes in Section~\ref{sec:lessons} were found on these traces, and the baselines received no such tuning. \sysname's thread\_finder uses a larger model than the baselines, which all use \texttt{gpt-6-luna}; the answer itself is written by the same model in every system. The baselines ran once each, for cost; \sysname's run-to-run spread ($\pm$0.005) suggests one run is representative for an LLM pipeline of this kind, but we did not measure theirs.

\textbf{External validity.} The programs are small and inject one fault at a time, whereas real systems often have several problems together and far more threads. All traces come from one machine and kernel and last 30\,s; our earlier prototype handled traces of up to 30 minutes, but not on this bench. All systems used one model family. The knowledge base includes 10 off-topic Django documents, which act as distractors, and no program documentation.

\textbf{Construct validity.} Verdicts are read from one line, so an answer that names the wrong thread still counts as a true positive. The bench includes an answer key with the threads, causes and numbers a correct answer should give; grading answers against it is future work.

\textbf{Open problems.} Several symptoms are hard to see in a kernel trace. Memory pressure needs page-fault and kernel-memory events that our kernel's LTTng lacks. Busy-waiting looks the same as computing. A trace does not record lock owners, so convoys and lock contention are inferred, and these caused the largest share of \sysname's persistent errors. Wake-ups through loopback networking end the critical path at ``network'' instead of the peer process. The critical path also stops at preemption: it shows that a thread on the path waited for a CPU, but not which thread held it, so on a priority inversion an answer can report the holder's 15\,ms wait for a CPU in every round without naming the real-time thread that preempted it. Beyond detection, production use raises questions this bench does not answer: traces from many hosts, user-space and kernel events together, and the cost of the thread search on systems with thousands of threads.
